% GENERATED FILE. Edit canonical lessons, then run npm run generate:books.
% canonical-source-hash: fnv1a64:5b42fdef530ed9b2
% canonical-lessons: KA-C36-magu, KA-C36-maga, KA-C36-magalu, KA-S131-vowel-sign-vocalic-r

\chapter{One Root, Three Endings}
\label{ch:ka-child-son-daughter}
\hlchaptermodality{36}

\begin{chapteropening}
\textbf{By the end of this chapter:} \emph{I can say child, son and daughter in Kannada, and show how one shared root with three different endings does the job Chapter 12 needed four separate words for.}

Line up \kn{ಮಗು}, \kn{ಮಗ} and \kn{ಮಗಳು} against Proto-Dravidian \emph{maka, and contrast gender-by-ending here with Chapter 12's four unrelated age-graded sibling words.}
\end{chapteropening}

\section[magu]{\kn{ಮಗು} (magu) — \textquotedblleft{}child,\textquotedblright{} before it picks a gender}

\label{lesson:KA-C36-magu}



\begin{quote}
\emph{Snēhita} and \emph{geḷeya} were two words for one idea. This word is the opposite kind of pair-maker: one word that the next two lessons will split in two.
\end{quote}

\subsection*{You'll want to know: \kn{ಮಗು}}
\begin{quote}
\textbf{\kn{ಮಗು}} — \emph{magu} — \textbf{baby, child} (no gender specified)
\end{quote}

\begin{quote}
\textbf{\kn{ಇದು} \kn{ನನ್ನ} \kn{ಮಗು}.} — \emph{idu nanna magu.} — \textquotedblleft{}This is my child.\textquotedblright{}
\end{quote}

\begin{cousinweb}
\textbf{\kn{ಮಗು}} goes back to Proto-Dravidian \textbf{*maka}, \textquotedblleft{}child\textquotedblright{} — a plain root with no gender marked on it at all. Hold that root in mind: the next lesson adds one ending to it and gets \textquotedblleft{}son\textquotedblright{}; the one after adds a different ending and gets \textquotedblleft{}daughter.\textquotedblright{} Chapter 12 solved the parallel problem — telling siblings apart — with four completely separate words. This root solves a smaller problem the opposite way: one shape, several endings.
\end{cousinweb}

\subsection*{Your turn}
Say these aloud:

\begin{itemize}
\raggedright
  \item \textquotedblleft{}idu nanna magu\textquotedblright{} — this is my child
  \item the root to hold onto — \textquotedblleft{}maka\textquotedblright{}
  \item friend, then child — \textquotedblleft{}snēhita … magu\textquotedblright{}
\end{itemize}

\begin{tcolorbox}[breakable,colback=teal!4,colframe=teal!35!black,title={Before you move on}]
Say \textquotedblleft{}this is my child.\textquotedblright{} (\emph{Idu nanna magu}.) Does \emph{magu} specify a gender? (\textbf{No.}) What is the Proto-Dravidian root underneath it, and what happens to it next? (\emph{\textbf{*maka}} — the next two lessons put different endings on it for \textquotedblleft{}son\textquotedblright{} and \textquotedblleft{}daughter.\textquotedblright{})
\end{tcolorbox}
\section[maga]{\kn{ಮಗ} (maga) — \textquotedblleft{}son,\textquotedblright{} the same root with one vowel doing the work}

\label{lesson:KA-C36-maga}



\begin{quote}
\emph{Maga}, \emph{magu} — hold the root steady and change almost nothing.
\end{quote}

\subsection*{You'll want to know: \kn{ಮಗ}}
\begin{quote}
\textbf{\kn{ಮಗ}} — \emph{maga} — \textbf{son}
\end{quote}

\begin{quote}
\textbf{\kn{ಇವನು} \kn{ನನ್ನ} \kn{ಮಗ}.} — \emph{ivanu nanna maga.} — \textquotedblleft{}He is my son.\textquotedblright{}
\end{quote}

Set beside \emph{magu}: \textbf{\kn{ಮಗು}} \emph{magu} $\to$ \textbf{\kn{ಮಗ}} \emph{maga}. The same three letters, minus the length on the final vowel — the plain, inherent \emph{-a} now standing for \textquotedblleft{}male.\textquotedblright{}

\begin{cousinweb}
\emph{Maga} descends from the same Proto-Dravidian \textbf{*maka} as \emph{magu}. The sisters carry a longer form for \textquotedblleft{}son\textquotedblright{} — reconstructed as \textbf{*magan} — and Kannada is the branch that wore the ending down to a single open vowel. Tamil keeps more of the shape: \textbf{\ta{மகன்}} (\emph{makaṉ}), \textquotedblleft{}son.\textquotedblright{}
\end{cousinweb}

\subsection*{Your turn}
Say these aloud:

\begin{itemize}
\raggedright
  \item \textquotedblleft{}ivanu nanna maga\textquotedblright{} — he is my son
  \item the near pair — \textquotedblleft{}magu … maga\textquotedblright{}
  \item Kannada's worn form, then Tamil's fuller one — \textquotedblleft{}maga … makaṉ\textquotedblright{}
\end{itemize}

\begin{tcolorbox}[breakable,colback=teal!4,colframe=teal!35!black,title={Before you move on}]
Say \textquotedblleft{}he is my son.\textquotedblright{} (\emph{Ivanu nanna maga}.) How does \emph{maga} differ from \emph{magu}? (\textbf{One vowel} — the plain \emph{-a} standing for \textquotedblleft{}male.\textquotedblright{}) What does Tamil's fuller form for \textquotedblleft{}son\textquotedblright{} look like? (\emph{\textbf{\ta{மகன்}}} \emph{makaṉ} — Kannada wore the ending down to just the vowel.)
\end{tcolorbox}
\section[magaḷu]{\kn{ಮಗಳು} (magaḷu) — \textquotedblleft{}daughter,\textquotedblright{} one more ending on the same root}

\label{lesson:KA-C36-magalu}



\begin{quote}
\emph{Magu} gave you the plain root, \emph{maga} the male ending. One more ending closes the set.
\end{quote}

\subsection*{You'll want to know: \kn{ಮಗಳು}}
\begin{quote}
\textbf{\kn{ಮಗಳು}} — \emph{magaḷu} — \textbf{daughter}
\end{quote}

\begin{quote}
\textbf{\kn{ಇವಳು} \kn{ನನ್ನ} \kn{ಮಗಳು}.} — \emph{ivaḷu nanna magaḷu.} — \textquotedblleft{}She is my daughter.\textquotedblright{}
\end{quote}

\begin{grammarlens}[title={Grammar Lens: gender by ending, not by a new word}]
Line the three up:

\noindent
\begin{tabularx}{\linewidth}{@{}>{\raggedright\arraybackslash}X>{\raggedright\arraybackslash}X@{}}
\toprule
\textbf{Kannada} & \textbf{means} \\
\midrule
\textbf{\kn{ಮಗು}} \emph{magu} & child, no gender \\
\textbf{\kn{ಮಗ}} \emph{maga} & son \\
\textbf{\kn{ಮಗಳು}} \emph{magaḷu} & daughter \\
\bottomrule
\end{tabularx}

One root, three endings. Chapter 12 solved a different split — \textbf{older} against \textbf{younger} sibling — the opposite way, with \textbf{four unrelated words}, \kn{ಅಣ್ಣ}/\kn{ತಮ್ಮ} and \kn{ಅಕ್ಕ}/\kn{ತಂಗಿ} sharing no root at all. Gender here rides on a suffix; age there needed a whole new word every time.
\end{grammarlens}

\begin{cousinweb}
\emph{Magaḷu} continues Proto-Dravidian \textbf{*makaḷ}, \textquotedblleft{}daughter.\textquotedblright{} Tamil \textbf{\ta{மகள்}} (\emph{makaḷ}) and Malayalam \textbf{\ml{മകൾ}} (\emph{makaḷ}) are close to identical to Kannada's own form — this branch of the family barely touched the word at all.
\end{cousinweb}

\subsection*{Your turn}
Say these aloud:

\begin{itemize}
\raggedright
  \item \textquotedblleft{}ivaḷu nanna magaḷu\textquotedblright{} — she is my daughter
  \item the set of three — \textquotedblleft{}magu … maga … magaḷu\textquotedblright{}
  \item one root three endings, then four words no root — \textquotedblleft{}maga, magaḷu … aṇṇa, tamma, akka, tangi\textquotedblright{}
  \item the sisters, near-identical — \textquotedblleft{}magaḷu … makaḷ … makaḷ\textquotedblright{}
\end{itemize}

\begin{tcolorbox}[breakable,colback=teal!4,colframe=teal!35!black,title={Before you move on}]
Say \textquotedblleft{}she is my daughter.\textquotedblright{} (\emph{Ivaḷu nanna magaḷu}.) Name the three words built on \emph{maka} and what each means. (\emph{Magu} — child; \emph{maga} — son; \emph{magaḷu} — daughter.) How does that differ from how Chapter 12 split siblings by age? (\textbf{Chapter 12 used four unrelated words}; this set uses \textbf{one root and different endings}.) How close is Tamil's word for daughter? (\textbf{Nearly identical} — \emph{makaḷ}.)
\end{tcolorbox}
\section[vowel sign ṛ]{\kn{◌ೃ} — one character, and the pair met the other way round}

\label{lesson:KA-S131-vowel-sign-vocalic-r}



\begin{quote}
Before the new one: \kn{ಓ} — does it stand alone or ride on a consonant, and what is its other shape?

One character this time, and today the halves arrive in the opposite order from every vowel so far.
\end{quote}

\subsection*{Script you'll notice: \kn{◌ೃ}}
\textbf{\kn{◌ೃ}} — the riding form of \textbf{\kn{ಋ}} (\emph{ṛ}).

Every vowel until now came standing first and riding second: \textbf{\kn{ಉ}} then \textbf{\kn{ು}}, \textbf{\kn{ಊ}} then \textbf{\kn{ೂ}}, \textbf{\kn{ಈ}} then \textbf{\kn{ೀ}}. This one is the reverse. You have had the standing letter since the chapter that taught \textbf{\kn{ಋ}}, in a word you say:

\begin{itemize}
\raggedright
  \item \textbf{\kn{ಋತು}} \emph{ṛtu} — season
\end{itemize}

What you have not seen is that vowel riding a consonant. It becomes a small hook tucked underneath:

\begin{itemize}
\raggedright
  \item \textbf{\kn{ಹ}} plus \textbf{\kn{◌ೃ}} gives \textbf{\kn{ಹೃ}} — \emph{hṛ}
\end{itemize}

Standing at the front of a word, look for \textbf{\kn{ಋ}}. Tucked under a consonant, look for \textbf{\kn{◌ೃ}}. Same vowel, two places it can sit.

\textbf{What it sounds like.} An \emph{r} doing a vowel's work — close to the \emph{ri} in a quick \emph{bring it}. Kannada took the sound and both shapes from Sanskrit, which is where the words needing it come from.

\subsection*{Writing: \kn{◌ೃ} — copy what you see}
Put your pen on \kn{◌ೃ} and follow its line. Copy the shape you can see — slowly, and larger than it is printed.

\begin{quote}
This book does not yet tell you \textbf{where to start the character or which way to travel}. That is a real thing, taught with real variation from school to school, and it is not written down here until it can be written down with a source. Copying what is in front of you needs no such source, and it is how the shape gets into your hand in the meantime.
\end{quote}

\subsection*{Your turn}
\begin{itemize}
\raggedright
  \item \emph{Look:} at these two and say which one is standing and which one is riding
\end{itemize}

\begin{quote}
\kn{ಋತು}  ·  \kn{ಹೃ}
\end{quote}

\begin{itemize}
\raggedright
  \item \emph{Trace it:} \kn{◌ೃ} three times under a \kn{ಹ}, saying \emph{hṛ} as you finish each one
  \item \emph{Look:} at \kn{ಋ} and \kn{◌ೃ} side by side, and say which one opens a word
\end{itemize}

\begin{tcolorbox}[breakable,colback=teal!4,colframe=teal!35!black,title={Before you move on}]
Which character is this — \kn{◌ೃ}? Which standing letter is it the other half of? (\textbf{\kn{ಋ}}.) Where does it sit on the page? (\textbf{Tucked under the consonant it rides.}) Name the word you already say that carries the standing one. (\emph{ṛtu}, season.)
\end{tcolorbox}
